\chapter{理论技术}
\vspace{7pt}

\section{日志概述}

服务器日志是系统在运行过程中产生的事件记录，通常包含时间、节点、组件、日志级别和日志正文等信息。日志既记录正常执行过程，也记录异常告警、错误堆栈和资源访问情况，是分析服务器运行状态的重要数据来源。与结构化监控指标相比，日志具有更丰富的语义信息，但也存在格式不统一、噪声较多和动态参数复杂等问题\cite{oliner2007supercomputers,xu2009detecting}。

服务器日志通常具有以下特点。

\begin{enumerate}
    \item \textbf{规模较大。} 服务器在持续运行过程中会产生大量日志，高并发服务、分布式部署和多节点集群会进一步增加日志规模。
    \item \textbf{格式多样。} 不同组件和程序模块输出的日志格式并不一致，同一系统内部也可能同时存在结构化、半结构化和自由文本日志。
    \item \textbf{时序性。} 日志按照时间顺序产生，事件之间的先后关系能够反映系统执行路径和故障传播过程。
    \item \textbf{采样不均匀。} 日志不是等间隔产生的数据。系统平稳运行时日志较为稀疏，异常发生时可能出现短时间高密度输出。
    \item \textbf{类别不平衡。} 正常日志通常占绝大多数，异常日志数量相对较少，模型容易偏向多数类。
\end{enumerate}

上述特点决定了日志异常检测需要先将原始文本转换为结构化表示，再结合统计、序列和时间特征进行建模。

\section{日志预处理}

原始日志不能直接用于模型训练。日志预处理的目标是消除无效字段和格式差异，保留与异常检测相关的信息，并为后续日志解析和特征构建提供规范输入。

\subsection{数据清洗}

日志数据清洗主要包括空值处理、重复记录过滤、异常字符清理和无关字段筛选。空值处理用于去除时间戳或日志正文缺失的记录；重复记录过滤用于减少重复输出对频次统计的干扰；异常字符清理用于统一空白符、控制字符和无意义符号；无关字段筛选用于降低输入维度。清洗过程不应改变异常标签和时间戳等监督学习所需字段。

\subsection{时间标准化}

日志具有时间属性，时间戳标准化是时序建模的基础。设原始日志集合为
\begin{equation}
    L=\{l_1,l_2,\cdots,l_n\},
\end{equation}
其中第 $i$ 条日志可表示为 $l_i=(t_i,c_i,m_i)$，$t_i$ 为时间戳，$c_i$ 为结构化字段，$m_i$ 为日志正文。经过时间排序后得到
\begin{equation}
    L'=sort(L,t_i).
\end{equation}
排序后的日志能够保留事件先后关系，并为时间窗口划分和相邻事件时间间隔计算提供依据。

\subsection{级别处理}

日志级别通常包括 INFO、WARN、ERROR、DEBUG 和 FATAL 等。日志级别能够反映事件严重程度，但不能简单等同于异常标签。例如，单条 ERROR 日志可能被自动恢复，而多个 INFO 事件的异常顺序也可能提示系统状态异常。因此，日志级别适合作为辅助特征，而不宜作为唯一判定规则。

\section{日志解析}

日志解析的目标是将非结构化日志正文转换为结构化事件。日志正文通常由固定模板和动态参数组成，其中固定模板描述事件语义，动态参数描述运行时变量。若不进行解析，同一事件会因不同 IP、端口、路径或编号被视为不同文本，导致模型输入空间膨胀。

\subsection{日志模板}

例如，以下两条日志的动态参数不同，但表示同一类连接失败事件。

\begin{verbatim}
2025-03-01 10:12:03 ERROR Failed to connect to 192.168.1.10:8080
2025-03-01 10:12:08 ERROR Failed to connect to 192.168.1.15:8080
\end{verbatim}

经过模板化处理后，可表示为
\begin{verbatim}
Failed to connect to <*>:<*>
\end{verbatim}
其中 \texttt{<*>} 表示被归一化的动态参数。设日志解析函数为 $P(\cdot)$，日志正文 $m_i$ 可映射为事件 ID：
\begin{equation}
    e_i=P(m_i).
\end{equation}

\subsection{解析方法}

常见日志解析方法包括规则方法、聚类方法、树结构方法和学习方法。规则方法依赖人工设计正则表达式，实现成本低，适合格式稳定的日志。聚类方法根据文本相似度或词项分布抽取模板，能够减少人工规则设计，但在大规模日志上计算成本较高。树结构方法通过分层匹配提升解析效率，适合在线日志处理。学习方法能够自动提取日志结构，但通常需要较多训练样本和额外计算资源\cite{zhu2019tools,du2016spell,he2017drain}。

本文面向毕业论文系统实现和实验验证，采用规则与正则表达式结合的解析方式。该方式实现过程明确，便于复现，并便于在后续实验中解释模板生成结果。

\section{特征工程}

日志解析后得到事件序列，但事件 ID 仍需转换为模型可处理的数值特征。本文从事件频次、事件顺序和时间间隔三个角度构建特征。

\subsection{事件频次}

事件频次特征用于统计时间窗口内不同事件模板的出现次数。已有系统日志研究表明，事件计数和模板分布能够反映系统在局部时间范围内的运行状态。设模板集合为 $V=\{v_1,v_2,\cdots,v_k\}$，第 $j$ 个时间窗口为 $W_j$，其频次向量表示为
\begin{equation}
    X_j=[x_{j1},x_{j2},\cdots,x_{jk}],
\end{equation}
其中 $x_{jp}$ 表示模板 $v_p$ 在窗口 $W_j$ 中出现的次数。该特征能够反映局部事件分布变化，适合捕捉异常爆发导致的频次突增。

\subsection{日志序列}

事件频次无法表达事件先后顺序。对于服务器日志而言，某些异常并不表现为单个事件异常，而表现为事件组合顺序异常。因此，本文将窗口内日志按照时间排序后构造事件序列\cite{du2017deeplog,meng2019loganomaly}
\begin{equation}
    S_j=\{e_{j1},e_{j2},\cdots,e_{js}\}.
\end{equation}
当窗口内事件数量不足固定长度时使用补齐标记；当事件数量超过固定长度时保留最近的连续片段，以保证输入维度一致。

\subsection{时间间隔}

日志事件之间的真实时间间隔能够描述系统状态变化速度。对于非规则采样序列，连续时间建模能够避免将不同实际时间跨度的事件间隔简单视为相同时间步。设连续事件 $e_{i-1}$ 和 $e_i$ 的时间戳分别为 $t_{i-1}$ 和 $t_i$，其时间间隔定义为\cite{rubanova2019latentode}
\begin{equation}
    \Delta t_i=t_i-t_{i-1}.
\end{equation}
异常爆发时，多个事件可能在短时间内密集出现；平稳运行时，相邻日志间隔可能较长。将 $\Delta t_i$ 显式输入模型，有助于模型区分事件顺序相同但发生节奏不同的样本。

\subsection{嵌入表示}

事件 ID 是离散变量，直接使用整数编号可能引入不存在的大小关系。为增强模型表达能力，本文使用嵌入层将事件 ID 映射为连续向量：
\begin{equation}
    h_i=Embedding(e_i), \quad h_i\in\mathbb{R}^{d}.
\end{equation}
嵌入向量与频次特征、时间间隔特征拼接后形成联合输入。

\section{异常检测}

异常检测的目标是识别与正常运行模式存在差异的样本。在服务器日志场景中，异常可以表现为单条严重错误日志，也可以表现为一段时间内事件分布、事件顺序或时间节奏的异常变化。

\subsection{异常类型}

本文将日志异常概括为三类。

\begin{enumerate}
    \item \textbf{点异常。} 单条日志本身具有异常含义，例如严重故障或系统崩溃信息。
    \item \textbf{上下文异常。} 日志事件本身并不罕见，但在特定时间或特定运行阶段出现时不符合系统状态。
    \item \textbf{序列异常。} 多条日志组合后的顺序、频率或时间间隔偏离正常模式。
\end{enumerate}

本文重点关注窗口级异常检测，即判断某一时间窗口或日志轨迹是否处于异常状态。

\subsection{类别不平衡}

服务器日志异常样本比例通常较低。若模型简单地将全部样本预测为正常，可能获得较高准确率，但无法识别真正异常。因此，模型评价不能只依赖 Accuracy，还需要结合 Precision、Recall 和 F1 值进行分析\cite{he2009imbalanced,saito2015precision}。

\subsection{评价指标}

在二分类检测中，Precision 和 Recall 定义为
\begin{equation}
    Precision=\frac{TP}{TP+FP}, \quad Recall=\frac{TP}{TP+FN}.
\end{equation}
F1 值为二者的调和平均数：
\begin{equation}
    F1=\frac{2\times Precision\times Recall}{Precision+Recall}.
\end{equation}
其中，TP 表示异常样本被正确识别为异常，FP 表示正常样本被误报为异常，FN 表示异常样本被漏报。本文模型输出连续异常评分，最终类别由阈值决定。为减少固定阈值带来的偏差，本文在验证集上选择使 F1 值较优的阈值。

\section{Liquid 架构}

Liquid 架构是一类连续时间递归神经网络。与普通离散时序模型不同，Liquid 模型通过状态演化方程描述隐藏状态随输入和时间变化的过程，适合处理不规则采样序列\cite{chen2018neuralode,hasani2021liquid}。

以 Liquid Time-Constant Networks 为例，其单个神经元状态可抽象表示为
\begin{equation}
    \frac{dx(t)}{dt}=-\left[\frac{1}{\tau}+f(I(t),x(t),\theta)\right]x(t)+f(I(t),x(t),\theta)A,
    \label{eq:ltc_core}
\end{equation}
其中，$x(t)$ 表示状态变量，$\tau$ 表示基础时间常数，$I(t)$ 表示输入特征，$A$ 表示静息状态，$f(\cdot)$ 表示可学习的非线性函数。由式 \ref{eq:ltc_core} 可见，模型状态变化不仅依赖当前输入，也受到时间常数影响。

在本文系统中，Liquid 架构与日志异常检测的结合方式如下。首先，日志事件 ID 经嵌入层映射为连续向量；其次，系统将时间间隔作为显式输入，使模型能够感知日志产生节奏；最后，模型根据窗口内事件序列输出异常评分，并通过验证集阈值完成异常判定。该机制使模型能够同时利用事件语义、事件顺序和非规则时间信息。

\section{本章小结}

本章围绕服务器日志异常检测所需理论和技术进行了说明。首先分析了服务器日志的规模、格式、时序、采样和类别分布特征，明确了原始日志需要经过清洗、解析和特征构建后才能用于模型训练。其次，介绍了日志模板、动态参数替换和事件 ID 映射等日志解析思想，说明了将非结构化日志转换为结构化事件序列的必要性。随后，本章从事件频次、事件顺序、时间间隔和嵌入表示四个方面说明了日志特征工程方法，并给出了异常检测任务中常用的评价指标。最后，本章介绍了 Liquid 架构的连续时间建模思想，说明了其与服务器日志非规则时间间隔特征之间的对应关系。

本章内容为后续系统设计提供理论基础。第三章将在本章基础上进一步给出系统总体架构、关键模块和算法流程；第四章将说明这些理论与设计如何落实为可运行的系统实现；第五章将通过实验验证相关设计的有效性。

需要强调的是，本文所采用的理论和技术并非相互独立。日志解析决定事件 ID 的稳定性，事件 ID 决定序列建模的输入空间，时间窗口决定样本粒度，时间间隔决定连续时间建模的信息来源，评价指标决定阈值选择方向。若其中任一环节设计不合理，都会影响最终异常检测效果。因此，后续章节将按照数据流顺序展开，使每个模块的设计目的、实现方式和实验验证形成对应关系。
